% ============================================================================
%  TUFT 全维统一场：Navier--Stokes 涡动力学与电磁洛伦兹耦合
%  —— 作用量泛函、拉格朗日密度与统一变分原理（论文片段）
%
%  编译方式：XeLaTeX（Overleaf 中请将编译器切换为 XeLaTeX）
%  依赖宏包：amsmath, amssymb, amsfonts, bm, physics, geometry（Overleaf 均已内置）
%
%  内容结构：
%    0. 公理基底（TUFT：空间光速螺旋流形假设）
%    1. 场变量定义
%    2. 作用量泛函与拉格朗日密度（field / matter / coupling）
%    3. 变分求导：欧拉--拉格朗日方程（NS 动量方程 + 洛伦兹力）
%    4. 不变量与守恒律（Noether 定理）
%    5. 自相似爆破解在 TUFT 作用量中的定位
%    6. 统一方程总表（TUFT 核心方程组）
%    7. 下一阶段方向
% ============================================================================
\documentclass[11pt]{ctexart}

\usepackage{amsmath,amssymb,amsfonts}
\usepackage{bm}
\usepackage{physics}
\usepackage[margin=2.5cm]{geometry}

% ---- 矢量与微分宏（矢量微分宏包） ----
\newcommand{\ve}[1]{\bm{#1}}                         % 粗体矢量
\newcommand{\Ltot}{\mathcal{L}}                      % 总拉格朗日密度
\newcommand{\Lf}{\mathcal{L}_{\mathrm{field}}}       % 场拉格朗日密度
\newcommand{\Lm}{\mathcal{L}_{\mathrm{matter}}}      % 物质拉格朗日密度
\newcommand{\Lc}{\mathcal{L}_{\mathrm{coupling}}}    % 耦合拉格朗日密度
\newcommand{\Vtopo}{V_{\mathrm{topo}}}               % 拓扑束缚势

\begin{document}

\title{\bfseries TUFT 全维统一场：Navier--Stokes 涡动力学与电磁洛伦兹耦合\\[6pt]
\large 作用量泛函、拉格朗日密度与统一变分原理}
\author{}
\date{}

\maketitle

\begin{abstract}
本文给出 TUFT（全维统一场，空间光速螺旋流形假设）的完整作用量表述。真空底层动力学客体为空间光速螺旋流场 $\ve{u}(\ve{x},t)$，其涡量 $\ve{\omega}=\nabla\times\ve{u}$ 被同构识别为磁场，标量势 $\phi$ 对应螺旋径向曲率伸缩场。总作用量 $S=S_{\mathrm{field}}+S_{\mathrm{matter}}+S_{\mathrm{coupling}}$ 由三项构成：真空螺旋流场动力学（还原不可压缩 Navier--Stokes 方程及其涡量输运方程）、拓扑螺旋结（物质）作用量、以及螺旋结与背景矢势的几何耦合项。对点粒子自由度变分，洛伦兹力 $\ve{F}=q\ve{E}+q\ve{v}\times\ve{B}$ 作为作用量的直接推论自动导出，而非额外假设。应用 Noether 定理得到能量、动量、角动量守恒与规范不变性。自相似爆破解在该框架中被定位为作用量有限的奇异有限时间解。
\end{abstract}

% ---------------------------------------------------------------------------
\section{公理基底（TUFT：空间光速螺旋流形假设）}

真空为 $3+1$ 维伪黎曼流形，其底层动力学客体是空间光速螺旋流场 $\ve{u}(\ve{x},t)$（流形上的切矢量场）。涡量定义为
\begin{equation}
\ve{\omega}=\nabla\times\ve{u},
\label{eq:vorticity}
\end{equation}
$\ve{\omega}$ 是螺旋环向挠率密度，同构于磁场 $\ve{B}$；标量势 $\phi$ 对应螺旋径向曲率伸缩场（与电场同源）。

机电对偶映射：
\begin{equation}
\ve{\omega}\leftrightarrow\ve{B},\qquad
\ve{u}\leftrightarrow\ve{v},\qquad
\text{拓扑荷 } q \leftrightarrow \text{螺旋缠绕数}.
\label{eq:duality}
\end{equation}

不可压缩约束（流体积守恒，涡管守恒）：
\begin{equation}
\nabla\cdot\ve{u}=0.
\label{eq:incompressible}
\end{equation}

% ---------------------------------------------------------------------------
\section{场变量定义（TUFT 场变量集合）}

\begin{itemize}
\item $\ve{u}(\ve{x},t)$：螺旋流切向速度场（Navier--Stokes 速度场）；
\item $\ve{\omega}=\nabla\times\ve{u}$：涡量/挠率场（等价磁场 $\ve{B}$）；
\item $\phi(\ve{x},t)$：标量曲率势（电场势，螺旋径向伸缩）；
\item $\psi(\ve{x},t)$：拓扑荷标量场（粒子/螺旋结分布，$\rho_q\propto\psi$）。
\end{itemize}

\noindent\textbf{矢量势引入（与电磁 4-势对齐）。}
定义 4-势 $A^{\mu}=(\phi,\ve{A})$，满足
\begin{equation}
\ve{\omega}=\nabla\times\ve{A},\qquad
\ve{E}=-\nabla\phi-\pd{\ve{A}}{t}.
\label{eq:4pot}
\end{equation}
TUFT 同构：$\ve{A}$ 是螺旋流矢势，$\ve{\omega}=\nabla\times\ve{A}$。

% ---------------------------------------------------------------------------
\section{作用量泛函与拉格朗日密度}

作用量 $S$ 为 $3+1$ 维时空积分，分解为三项：
\begin{equation}
S=S_{\mathrm{field}}+S_{\mathrm{matter}}+S_{\mathrm{coupling}},
\label{eq:Sdecomp}
\end{equation}
其中
\begin{itemize}
\item $S_{\mathrm{field}}$：真空螺旋流场自身动力学作用量（含 NS 涡动能、曲率势能、粘性正则项）；
\item $S_{\mathrm{matter}}$：拓扑螺旋结（物质/电荷）的粒子作用量；
\item $S_{\mathrm{coupling}}$：物质螺旋结与背景真空螺旋场的几何耦合项（洛伦兹力的来源）。
\end{itemize}
\begin{equation}
S=\int_{t_1}^{t_2}\!\int_{\Omega}
\Ltot\!\left(\ve{u},\pd{\ve{u}}{t},\nabla\ve{u},\phi,\nabla\phi,\psi\right)\,d^{3}x\,dt,
\label{eq:Saction}
\end{equation}
其中 $\Ltot$ 为拉格朗日密度。

\subsection{场拉格朗日密度 $\Lf$}

真空螺旋场包含：涡旋动能项、径向曲率势能项、粘性耗散正则项、不可压缩约束拉格朗日乘子：
\begin{equation}
\Lf
=\frac{1}{2}\rho|\ve{u}|^{2}
-\frac{1}{2}\varepsilon|\nabla\phi|^{2}
+\frac{1}{2\mu}|\ve{\omega}|^{2}
-\lambda(\nabla\cdot\ve{u})
-\frac{1}{2}\rho\nu|\nabla\ve{u}|^{2}.
\label{eq:Lfield}
\end{equation}

逐项物理注释（螺旋几何解释）：
\begin{enumerate}
\item $\frac{1}{2}\rho|\ve{u}|^{2}$：螺旋流平移动能（NS 流体动能，对应 LC 电路磁场动能 $\tfrac12 L\dot{q}^{2}$）；
\item $-\frac{1}{2}\varepsilon|\nabla\phi|^{2}$：径向曲率伸缩势能（电场势能，对偶电容储能 $\tfrac12 q^{2}/C$）；
\item $\frac{1}{2\mu}|\ve{\omega}|^{2}$：涡旋挠率能；$\ve{\omega}=\nabla\times\ve{u}$ 即空间螺旋环向扭转储存的能量，电磁对应 $\frac{B^{2}}{2\mu_0}$；
\item $-\lambda(\nabla\cdot\ve{u})$：拉格朗日乘子，强制不可压缩条件 $\nabla\cdot\ve{u}=0$；
\item $-\frac{1}{2}\rho\nu|\nabla\ve{u}|^{2}$：粘性正则项，对应 NS 中的 $\nu\Delta\ve{u}$，代表螺旋曲面形变的耗散。注意：粘性是非守恒耗散项，在保守拉格朗日体系中只能作为低阶扰动/正则项；在 OpenAI 爆破构造中，粘性相对于非线性涡拉伸是微小扰动。
\end{enumerate}

\noindent\textbf{保守无粘子情况}：令 $\nu\to0$，得到 Euler 螺旋流场拉格朗日，对应无粘 NS 方程。

\subsection{物质拓扑螺旋结拉格朗日 $\Lm$}

拓扑荷（电荷/粒子）作为局域束缚螺旋孤子：
\begin{equation}
\Lm=\psi\left(\frac{1}{2}m|\dot{\ve{r}}|^{2}-\Vtopo(\psi)\right),
\label{eq:Lmatter}
\end{equation}
其中 $\Vtopo(\psi)$ 为拓扑束缚势，维持螺旋结不扩散，对应粒子质量的拓扑起源；$\psi$ 为孤子密度场。

点粒子极限 $\psi\to\delta(\ve{x}-\ve{r}(t))$，退化为离散质点：
\begin{equation}
L_{\mathrm{particle}}=\frac{1}{2}m\dot{\ve{r}}^{2}.
\label{eq:Lparticle}
\end{equation}

\subsection{几何耦合项 $\Lc$（洛伦兹力本源耦合）}

TUFT 核心耦合项，完全对标电磁最小耦合。几何意义：拓扑螺旋结与背景真空螺旋矢势的缠绕耦合：
\begin{equation}
\Lc=q\ve{A}\cdot\dot{\ve{r}}-q\phi.
\label{eq:Lcoupling}
\end{equation}
其中 $\ve{A}$ 为背景螺旋矢势，$q$ 为螺旋结的拓扑缠绕数；该项描述粒子螺旋与背景涡旋螺旋之间的拓扑链接数相互作用。

\subsection{完整总拉格朗日密度}

\begin{equation}\label{eq:totalL}
\begin{aligned}
\Ltot
={}&\underbrace{\frac{1}{2}\rho|\ve{u}|^{2}-\frac{1}{2}\varepsilon|\nabla\phi|^{2}+\frac{1}{2\mu}|\nabla\times\ve{u}|^{2}
+\lambda(\nabla\cdot\ve{u})-\frac{1}{2}\rho\nu|\nabla\ve{u}|^{2}}_{\Lf}\\
&+\underbrace{\psi\left(\frac{1}{2}m|\dot{\ve{r}}|^{2}-\Vtopo(\psi)\right)}_{\Lm}
+\underbrace{q\ve{A}\cdot\dot{\ve{r}}-q\phi}_{\Lc},
\end{aligned}
\end{equation}
其中 $\ve{\omega}=\nabla\times\ve{u}$，矢势 $\ve{A}$ 满足 $\nabla\times\ve{A}=\ve{\omega}$。

% ---------------------------------------------------------------------------
\section{变分求导：欧拉--拉格朗日方程（场方程 + 粒子运动方程）}

\subsection{场变量的欧拉--拉格朗日方程（$\to$ Navier--Stokes 方程）}

对连续场变量 $\ve{u}(\ve{x},t)$，欧拉--拉格朗日方程为
\begin{equation}
\frac{\partial\Ltot}{\partial\ve{u}}
-\partial_{\mu}\frac{\partial\Ltot}{\partial(\partial_{\mu}\ve{u})}=0,\qquad
\partial_{\mu}=(\partial_{t},\nabla).
\label{eq:ELfield}
\end{equation}

对 $\ve{u}$ 变分，代入 $\Ltot$ 分步求导：
\begin{equation}
\frac{\partial\Ltot}{\partial\ve{u}}=\rho\ve{u},
\label{eq:du}
\end{equation}
\begin{equation}
\frac{\partial\Ltot}{\partial(\partial_{t}\ve{u})}=0,\qquad
\frac{\partial\Ltot}{\partial(\partial_{i}u_{j})}
=\frac{1}{\mu}\varepsilon_{jkl}\,\omega_{l}\,\varepsilon_{kmi}\,\delta_{im}
-\rho\nu\,\partial_{i}u_{j}+\lambda\delta_{ij}.
\label{eq:dgradu}
\end{equation}
\footnote{【整理注】此式按原文手稿转录。展开后该项应化为挠率能的变分结构
$\frac{1}{\mu}\varepsilon_{jkl}\,\partial_{k}\omega_{l}=\frac{1}{\mu}(\nabla\times\ve{\omega})_{j}$，
即 $\frac{1}{\mu}\nabla\times(\nabla\times\ve{u})$ 的分量形式；结合 $\nabla\cdot\ve{u}=0$
给出 $\frac{1}{\mu}\Delta\ve{u}$ 型项及可并入压强 $p$ 的项。投稿前建议核对缩并与符号约定。}

执行分部积分并整理，最终得到动量方程：
\begin{equation}
\rho\left(\partial_{t}\ve{u}+(\ve{u}\cdot\nabla)\ve{u}\right)
=-\nabla p+\rho\nu\Delta\ve{u}+\ve{f}_{\mathrm{couple}},
\label{eq:NS}
\end{equation}
其中 $\ve{f}_{\mathrm{couple}}$ 是拓扑荷螺旋结施加给流场的外力；压强 $p$ 吸收拉格朗日乘子 $\lambda$ 与曲率势贡献。\emph{这正是三维不可压缩 Navier--Stokes 方程。}

对动量方程取旋度，直接得到涡量输运方程：
\begin{equation}
\partial_{t}\ve{\omega}+(\ve{u}\cdot\nabla)\ve{\omega}
=(\ve{\omega}\cdot\nabla)\ve{u}+\nu\Delta\ve{\omega}+\nabla\times\ve{f}_{\mathrm{couple}}.
\label{eq:vorticity-transport}
\end{equation}

\noindent\textbf{几何含义}：从统一作用量变分，自动还原 NS 涡动力学。

\subsection{点粒子变分：导出洛伦兹力（核心结论）}

对点粒子自由度 $\ve{r}(t)$ 变分。离散粒子拉格朗日量：
\begin{equation}
L(\ve{r},\dot{\ve{r}})=\frac{1}{2}m\dot{\ve{r}}^{2}+q\ve{A}(\ve{r},t)\cdot\dot{\ve{r}}-q\phi(\ve{r},t).
\label{eq:Lpoint}
\end{equation}
欧拉--拉格朗日方程：
\begin{equation}
\dv{}{t}\!\left(\pd{L}{\dot{\ve{r}}}\right)-\pd{L}{\ve{r}}=0.
\label{eq:ELparticle}
\end{equation}

\textbf{第一步：求广义动量。}
\begin{equation}
\ve{P}=\pd{L}{\dot{\ve{r}}}=m\dot{\ve{r}}+q\ve{A}.
\label{eq:momentum}
\end{equation}

\textbf{第二步：求时间全导数。}
\begin{equation}
\dv{\ve{P}}{t}=m\ddot{\ve{r}}+q\dv{\ve{A}}{t},
\label{eq:dPdt}
\end{equation}
全导数链式法则：
\begin{equation}
\dv{\ve{A}}{t}=\partial_{t}\ve{A}+(\dot{\ve{r}}\cdot\nabla)\ve{A}.
\label{eq:chain}
\end{equation}

\textbf{第三步：求空间梯度项。}
\begin{equation}
\pd{L}{\ve{r}}=q\nabla(\ve{A}\cdot\dot{\ve{r}})-q\nabla\phi.
\label{eq:dLdr}
\end{equation}

代入欧拉--拉格朗日方程：
\begin{equation}
m\ddot{\ve{r}}+q\left(\partial_{t}\ve{A}+(\dot{\ve{r}}\cdot\nabla)\ve{A}\right)
-q\nabla(\ve{A}\cdot\dot{\ve{r}})+q\nabla\phi=0.
\label{eq:ELexpand}
\end{equation}

矢量恒等式（令 $\ve{v}=\dot{\ve{r}}$）：
\begin{equation}
\nabla(\ve{A}\cdot\ve{v})=(\ve{v}\cdot\nabla)\ve{A}+\ve{v}\times(\nabla\times\ve{A}),
\label{eq:vectorid}
\end{equation}
代入化简：
\begin{equation}
m\ddot{\ve{r}}=q\left(-\nabla\phi-\partial_{t}\ve{A}\right)-q\,\dot{\ve{r}}\times(\nabla\times\ve{A}).
\label{eq:Lforce-step}
\end{equation}

定义场同构：
\begin{equation}
\ve{E}=-\nabla\phi-\partial_{t}\ve{A},\qquad
\ve{B}=\nabla\times\ve{A}=\ve{\omega},
\label{eq:EB}
\end{equation}
得到洛伦兹力的完整表达式：
\begin{equation}
\boxed{\;\ve{F}=m\ddot{\ve{r}}=q\ve{E}+q\ve{v}\times\ve{B}\;}
\label{eq:Lorentz}
\end{equation}

\noindent\textbf{结论}：洛伦兹力是 TUFT 统一作用量变分后的直接推论，而非额外假设。
\begin{itemize}
\item $q\ve{E}$：径向曲率伸缩场的梯度力（可做功，螺旋拉伸势能）；
\item $q\ve{v}\times\ve{B}$：环向挠率涡旋横向耦合项（不做功，螺旋曲面切向偏转，与 NS 涡升力同源）。
\end{itemize}

% ---------------------------------------------------------------------------
\section{不变量与守恒律（Noether 定理）}

对作用量 $S$ 应用诺特定理：
\begin{enumerate}
\item \textbf{时间平移不变}：总能量守恒（螺旋场动能 + 曲率势能 + 孤子物质能量；OpenAI NS 爆破：全局能量守恒，局部梯度奇点完全兼容）；
\item \textbf{空间平移不变}：总动量守恒（场动量 + 粒子动量，实物 $\leftrightarrow$ 真空螺旋场交换动量）；
\item \textbf{空间旋转不变}：总角动量守恒（螺旋涡旋角动量）；
\item \textbf{规范不变性}：矢势规范变换
\begin{equation}
\ve{A}\mapsto\ve{A}+\nabla\chi,\qquad
\phi\mapsto\phi-\partial_{t}\chi,
\label{eq:gauge}
\end{equation}
作用量不变；等价地，涡旋拓扑链接数不变。
\end{enumerate}

\noindent\textbf{物理注释}：规范不变性的本质是螺旋流形的微分同胚冗余；矢势 $\ve{A}$ 不是直接可观测的几何量，可观测的是曲率 $\ve{E}$ 与挠率涡量 $\ve{B}=\ve{\omega}$。

% ---------------------------------------------------------------------------
\section{自相似爆破解在 TUFT 作用量中的定位}

回到 OpenAI NS 爆破：在作用量框架下，爆破解是场方程的一类奇异有限时间解。在 $\tau=T^{*}-t\to0$ 极限下，场的局部变分二阶导数趋于无穷（曲率/挠率发散），但作用量泛函本身依然有限，因为奇异区域被限制在紧集内，积分收敛。

\begin{itemize}
\item \textbf{无粘极限 $\nu\to0$}：爆破由非线性涡拉伸项 $(\ve{\omega}\cdot\nabla)\ve{u}$ 主导；
\item \textbf{微小粘性}：作为正则扰动，不消除爆破轮廓（OpenAI 核心论证）。
\end{itemize}

% ---------------------------------------------------------------------------
\section{统一方程总表（TUFT 核心方程组）}

\begin{equation}\label{eq:summary}
\boxed{\begin{aligned}
&S=\int dt\,d^{3}x\ \Ltot,\\
&\Ltot=\frac{1}{2}\rho|\ve{u}|^{2}-\frac{1}{2}\varepsilon|\nabla\phi|^{2}
+\frac{1}{2\mu}|\nabla\times\ve{u}|^{2}
+\lambda(\nabla\cdot\ve{u})-\frac{1}{2}\rho\nu|\nabla\ve{u}|^{2}\\
&\qquad\qquad+\psi\left(\frac{1}{2}m|\dot{\ve{r}}|^{2}-\Vtopo(\psi)\right)
+q\ve{A}\cdot\dot{\ve{r}}-q\phi,\\
&\ve{\omega}=\nabla\times\ve{u},\qquad \ve{B}=\ve{\omega},\qquad
\ve{E}=-\nabla\phi-\partial_{t}\ve{A},\\
&\rho\left(\partial_{t}\ve{u}+(\ve{u}\cdot\nabla)\ve{u}\right)
=-\nabla p+\rho\nu\Delta\ve{u}+\ve{f}_{\mathrm{couple}}
\quad(\text{NS 动量方程}),\\
&\nabla\cdot\ve{u}=0\quad(\text{不可压缩约束}),\\
&\partial_{t}\ve{\omega}+(\ve{u}\cdot\nabla)\ve{\omega}
=(\ve{\omega}\cdot\nabla)\ve{u}+\nu\Delta\ve{\omega}+\nabla\times\ve{f}_{\mathrm{couple}}
\quad(\text{涡量方程}),\\
&m\ddot{\ve{r}}=q\ve{E}+q\ve{v}\times\ve{B}
\quad(\text{洛伦兹力，粒子变分结果}).
\end{aligned}}
\end{equation}

% ---------------------------------------------------------------------------
\section{下一阶段方向}

\begin{enumerate}
\item 将本 TUFT 作用量、欧拉--拉格朗日方程与诺特守恒律整理为完整可编译 LaTeX 论文片段（本稿即为此工作，可直接导入 Overleaf）；
\item 引入挠率张量（Cartan 几何），将 TUFT 从矢量场版本升级为黎曼--嘉当流形版本，把引力一并纳入同一作用量；
\item 数值变分仿真：用 Python 离散化作用量、最小化作用量，重现自相似螺旋涡爆破轮廓；
\item 对作用量做规范固定，推导哈密顿形式，构造泊松括号，做量子化版本 TUFT。
\end{enumerate}

% ============================================================================
% 【整理注】如需补充参考文献，可在此处添加，例如：
% \begin{thebibliography}{9}
% \bibitem{nsblowup}（请补充 OpenAI NS 爆破论文的完整引用信息）
% \end{thebibliography}
% ============================================================================

\end{document}
